\chapter{INTRODUCTION}

\section{Background}
\label{intro}
Low back pain (LBP) is a major health problem worldwide. Studies report that 60\% to 80\% of individuals experience it at some point during their lifetime \cite{Freeman2010}. The Lumbar Multifidus Muscle (LMM), which belongs to the deep spinal muscle group, plays an important role in maintaining posture and stabilizing the lower back. Clinical and imaging studies have shown that chronic LBP is often associated with multifidus muscle atrophy and fatty infiltration \cite{Vohra2016}. MRI-based observations also suggest that degeneration is often prominent near the L5 vertebral level, which is highly relevant for lumbar stability \cite{Vohra2016}.

CT scans provide good visualization of the LMM, but the use of ionizing radiation limits their routine use \cite{Brenner2022}. MRI is non-invasive and detailed, but it is expensive and may not be suitable for patients with metal implants, claustrophobia, or limited healthcare coverage. Ultrasound imaging offers a practical alternative because it is safe, inexpensive, portable, and suitable for repeated assessment.

However, ultrasound interpretation is not straightforward. Image quality varies, speckle noise is common, and the distinction between healthy muscle and fatty infiltration can be subtle. Manual analysis also depends on the skill and consistency of the examiner. A machine learning based approach can reduce this dependency by extracting consistent image features and producing repeatable assessments. Texture analysis and pixel-level study of ultrasound images have already shown promise for measuring muscle composition \cite{Pincheira2024}. This project builds on that direction by developing an automated pipeline for LMM localization and Heckmatt grade classification.

\section{Motivation}
\label{motivation}
Accurate characterization of the multifidus muscle is useful for understanding chronic back pain and for monitoring rehabilitation. MRI and CT provide high-quality images, but their cost, availability, and practical limitations reduce their usefulness for regular screening. Ultrasound is more accessible, but separating muscle tissue from intramuscular fat remains difficult, especially when image quality differs across machines and operators \cite{Sions2017}.

Deep learning can help make ultrasound-based assessment more objective. A model that first identifies the LMM region and then classifies muscle quality can reduce manual effort and make the assessment more consistent. Such a system would be especially useful in clinical settings where physicians need quick and repeatable support rather than a complex research workflow.

\section{Problem Definition}
\label{problem}
The problem addressed in this project is the automated grading of Lumbar Multifidus Muscle quality from B-mode ultrasound images. The main objective is to assign the image to one of the four Heckmatt grades by analysing the echo intensity and texture patterns within the muscle region.

\section{Thesis Organization}
This thesis is organized as follows. Chapter 2 reviews related work on lumbar multifidus imaging, ultrasound analysis, segmentation, classification, and imbalanced learning. Chapter 3 describes the proposed methodology for ROI extraction and Heckmatt grading. Chapter 4 presents the evaluation metrics used for segmentation and classification. Chapter 5 summarizes the experimental work and implementation details. Chapter 6 discusses the major results and observations. Chapter 7 concludes the thesis and outlines future directions.
